% Werke X/1 (1917), native188–190, printed181–183.
\begingroup
\setlength{\headheight}{25pt}
\SourceLeaf{188}\NoSourceHead
\SourceGroup{III}
\HistoricTitle{[DIFFERENTIAL- UND FUNKTIONALGLEICHUNGEN.]}
\HistoricCitation{[Zwei Zettel in Ff, Kapsel 46a.]}
\NumberedItem{1}
\noindent [Erster Zettel]
\par\indent Ponendo
\[\begin{aligned}
y&=x^{-1}+\frac14x^{-3}+\frac14.\frac9{16}x^{-5}+\frac14.\frac9{16}.\frac{25}{36}x^{-7}\text{ etc.}\\
&\hspace{17em}(=\varphi x
\end{aligned}\]
\noindent eruitur
\[xy+(3xx-1)\frac{dy}{dx}+x(xx-1)\frac{d^2y}{dx^2}=0.\]
\noindent Iam ponendo \(x=\frac12(t+t^{-1})\) erit
\[\begin{aligned}
\frac{dt}{dx}&=\frac2{1-t^{-2}}\\
\frac{dy}{dx}&=\frac{dy}{dt}.\frac2{1-t^{-2}}\\
\frac{d^2y}{dx^2}&=\left(\frac2{1-t^{-2}}\right)^2\frac{d^2y}{dt^2}
-\frac1{t^3}\left(\frac2{1-t^{-2}}\right)^3\frac{dy}{dt}.
\end{aligned}\]
\noindent Hinc
{\small
\[\begin{aligned}
&\frac{tt+1}{2t}y+\frac{3t^4+2tt+3}{2tt-2}.\frac{dy}{dt}\\
&\quad+\frac{tt+1}{2t}\left(\frac{tt-1}{2t}\right)^2
\left\{\left(\frac{2tt}{tt-1}\right)^2\frac{d^2y}{dt^2}
-\left(\frac{2t}{tt-1}\right)^3\frac{dy}{dt}\right\}
\end{aligned}\]}
\noindent sive
\[\frac{tt+1}{2t}y+\frac{3t^4+1}{2tt-2}.\frac{dy}{dt}+\frac{t(tt+1)}2\frac{d^2y}{dt^2}=0\]
\noindent sive
\[(t^4-1)y+t(3t^4+1)\frac{dy}{dt}+tt(t^4-1)\frac{d^2y}{dt^2}=0.\]
\noindent Hinc facile deducitur
\clearpage
\SourceLeaf{189}\ArithmeticHead{182}{L}{ANALYSIS.\quad NACHLASS.}
\[y=2\left(t^{-1}+\frac14t^{-5}+\frac14.\frac9{16}t^{-9}+\frac14.\frac9{16}.\frac{25}{36}t^{-13}+\text{etc.}\right)\]
\noindent sive
\[\varphi\frac12\left(t+\frac1t\right)=2t\varphi tt.\]
\noindent Ponendo \(x=\frac1p\) erit
\[y-p(1+pp)\frac{dy}{dp}+pp(1-pp)\frac{d^2y}{dp^2}=0.\]
\ClosingRule
\NumberedItem{2}
\noindent [Zweiter Zettel]
\par\indent Wenn
\[\varphi x=x^{-1}+\frac14x^{-3}+\frac{1.9}{4.16}x^{-5}+\text{etc.},\]
\noindent so ist bewiesen, dass
\[2x\varphi xx=\varphi\frac{xx+1}{2x}\quad .\,.\,.\,.\,.\,.\,.\,.\,.\,.\ (1).\]
\par\indent Es sei
\[\frac{\sqrt{xx-1}}{\varphi\surd\frac{xx}{xx-1}}=\Psi x,\]
\noindent woraus, \(\surd\frac{xx}{xx-1}=u\) gesetzt, folgt
\[\varphi u=\frac1{\sqrt{uu-1}.\Psi\surd\frac{uu}{uu-1}}.\]
\noindent Dann gibt die Functionalgleichung (1)
\[\frac{2x}{\sqrt{x^4-1}.\Psi\surd\frac{x^4}{x^4-1}}
=\frac1{\frac{xx-1}{2x}\Psi\frac{xx+1}{xx-1}},\]
\noindent mithin
\[\Psi\surd\frac{x^4}{x^4-1}=\surd\frac{xx-1}{xx+1}.\Psi\frac{xx+1}{xx-1}.\]
\noindent Man setze \(\frac{xx-1}{xx+1}=tt\), so wird \(xx=\frac{1+tt}{1-tt}\) und
\[\Psi\frac{1+tt}{2t}=\frac1t\Psi tt.\]
\par\indent Nun ist aber auch
\[M\frac{1+tt}{2t}=\frac1t Mtt.\]
\clearpage
\SourceLeaf{190}\ArithmeticHead{183}{R}{\scriptsize DIFFERENTIAL- UND FUNKTIONALGLEICHUNGEN.}
\noindent Macht man daher allgemein
\[\frac{\Psi z}{Mz}=Fz,\]
\noindent so wird
\[F\frac{1+tt}{2t}=Ftt,\]
\noindent woraus man leicht schliesst, dass \(Fz\) eine Konstante sein müsse.
\par\indent Nun ist für \(u=\infty\), \(\varphi u=\frac1u\); also \(\Psi\surd\frac{uu}{uu-1}\) oder \(\Psi1=1\).
Es ist aber auch \(M1=1\), folglich die Konstante \(=1\) und daher
\[Mx=\frac{\sqrt{xx-1}}{\varphi\surd\frac{xx}{xx-1}}.\]
\noindent Also [ist] \(\frac1{Mx}\) der von \(V\) freie Theil des Ausdrucks
\[\frac1{\surd xx-1.\surd\left(\frac{xx}{xx-1}-\cos V^2\right)}\]
\noindent oder auch der von \(V\) freie Theil bei
{\small\[\frac1{\surd(xx-\cos V^2(xx-1))}
=\frac1{\surd(xx\sin V^2+\cos V^2)}
=\frac1{\surd(1+(xx-1)\sin V^2)},\]}
\noindent und folglich
\[\frac1{M\surd(1+y)}=\frac1{\surd(1+y\sin V^2)}
\quad\begin{gathered}\text{(abiectis partibus angulum }V\\\text{implicantibus)}\end{gathered}\]
\noindent Übrigens ist es gleichgültig, ob man \(\sin V^2\) oder \(\cos V^2\) schreibt.
\[\varphi x=1+\frac14x^2+\frac{1.9}{4.16}x^4\text{ etc.}\]
\[\varphi s=\frac1{Mc}=\frac{\operatorname{tang}}{sM\frac1c}\]
\[\begin{array}{rr}
0&0\\
30^0&0{,}728395\\
45^0&0{,}847213\\
60^0&0{,}931808\\
90^0&1.
\end{array}\]
\ClosingRule
\clearpage
\endgroup

